% =============================================================================
%  Irradiation Studies and Design Optimization of the ATLAS Tile Calorimeter
%  for the High-Luminosity LHC
%
%  Author:  Róbert Astaloš, on behalf of the ATLAS Tile Calorimeter System
%  Style:   tikzposter, A0 portrait, ATLAS colour scheme
% =============================================================================

\documentclass[a0paper,portrait, blockverticalspace=4em, colspace=2em]{tikzposter}
\tikzposterlatexaffectionproofoff
\usepackage{graphicx}
\usepackage{amsmath}
\usepackage{amssymb}
\usepackage{hyperref}
\usepackage{multicol}
\usetikzlibrary{calc}

% ---------------------------------------------------------------------------
%  Colours
% ---------------------------------------------------------------------------
\definecolor{maincolor}{HTML}{9e2b2f}
\definecolor{secondarycolor}{RGB}{158, 43, 47}
\definecolor{lightred}{RGB}{237, 199, 201}

\definecolor{tid}{RGB}{254,128,2}   
\definecolor{niel}{RGB}{0,140,0}    
\definecolor{see}{RGB}{190,0,62}    
\definecolor{sel}{RGB}{0,90,170}    

% ===========================================================================
%  LAYOUT PARAMETERS
% ===========================================================================
\newcommand{\looseitems}{\small\begin{itemize}\setlength{\itemsep}{0.35em}}
\newcommand{\tightitems}{\small\begin{itemize}\setlength{\itemsep}{0.2em}}
\newcommand{\captiontext}[1]{%
    {\footnotesize\renewcommand{\baselinestretch}{0.90}\selectfont#1}}

% ---------------------------------------------------------------------------
%  Theme
% ---------------------------------------------------------------------------
\usetheme{Default}
\definecolorstyle{atlascolors}{
    \colorlet{backgroundcolor}{white}
    \colorlet{titlefgcolor}{white}
    \colorlet{titlebgcolor}{maincolor}
    \colorlet{blocktitlefgcolor}{white}
    \colorlet{blocktitlebgcolor}{maincolor}
    \colorlet{blockbodyfgcolor}{black}
    \colorlet{blockbodybgcolor}{lightred!25}
}{}
\usecolorstyle{atlascolors}

\definetitlestyle{sampletitle}{width=760mm, roundedcorners=20, linewidth=2pt,
  innersep=10pt, titletotopverticalspace=6mm, titletoblockverticalspace=8mm}{%
  \begin{scope}[line width=\titlelinewidth, rounded corners=\titleroundedcorners]
    \draw[color=blocktitlebgcolor, fill=titlebgcolor]
      (\titleposleft,\titleposbottom) rectangle (\titleposright,\titlepostop);
  \end{scope}
  \node[anchor=east, fill=white, rounded corners=10pt, inner sep=10pt, xshift=-22mm]
    at ($(\titleposright,\titlepostop)!0.5!(\titleposright,\titleposbottom)$)
    {
\includegraphics[height=8.6cm]{logos/atlas_transparent.png}};}
\usetitlestyle{sampletitle}

\title{\parbox{0.74\linewidth}{\centering\huge
    Irradiation Studies and Design Optimization of the\\[2mm]ATLAS Tile Calorimeter for the High-Luminosity LHC}}
\author{Róbert Astaloš \quad on behalf of the ATLAS Tile Calorimeter System}
\institute{Comenius University Bratislava}
\date{}

\setlength{\columnsep}{0.5cm}

\begin{document}
\maketitle

\begin{columns}

% ===========================================================================
%  COLUMN 1 
% ===========================================================================
\column{0.333}

\block{The ATLAS Tile Calorimeter \& the HL-LHC Upgrade}{
    \looseitems
        \item The ATLAS Tile Calorimeter (TileCal) provides crucial measurements of hadronic jets and missing transverse energy in the central region ($|\eta|<1.7$).
        \item For the \textbf{High-Luminosity LHC}, the on-detector electronics are being fully replaced to support un-triggered digital read-out of all channels at \textbf{40\,MHz} with improved timing resolution.
        \item This new electronics architecture must survive an unprecedented radiation environment for over a decade of operation without intervention.
        \item Consequently, rigorous qualification methodologies are applied to ensure reliability against \textbf{\textcolor{tid}{Total Ionizing Dose (TID)}}, \textbf{\textcolor{niel}{Non-Ionizing Energy Loss (NIEL)}}, and \textbf{\textcolor{see}{Single Event Effects (SEE)}}.
    \end{itemize}
}

\block{Data Read-Out Architecture}{
    \looseitems
        \item The upgraded system relies on the Daughterboard (DB) acting as the central hub within each Minidrawer.
        \item It aggregates digitized signals from the Front-End Boards and transmits them off-detector via redundant multi-Gbps optical links.
        \item Concurrently, the DB receives the LHC clock, configuration commands, and control signals from the Tile PreProcessor (TilePPr).
        \item Maintaining the synchronization of this read-out chain under high radiation is the primary focus of the qualification campaigns.
        \item Continuous un-triggered data transmission demands extreme reliability; a failure in the DB severs the link for up to 12 PMTs.
    \end{itemize}
}

\block{Radiation Environment \& Criteria}{
    \looseitems
        \item Radiation levels vary significantly depending on the position within the TileCal geometry. Criteria are defined using simulations for the \emph{most-exposed} regions.
        \item Qualification levels incorporate rigorous safety factors (SF) to account for:
        \begin{itemize}
            \normalsize
            \item \textbf{1.5} for pure simulation uncertainty.
            \item \textbf{3.0 -- 5.0} for low-dose-rate effects and lot-to-lot component variation across the production life cycle.
            \item \textbf{1.3} for non-monoenergetic beam equivalence.
        \end{itemize}
        \item The final Daughterboard qualification targets reach \textbf{\textcolor{tid}{TID $\approx$ 108\,Gy}}, \textbf{\textcolor{niel}{NIEL $\approx 13.16\times10^{12}$\,n/cm$^2$}}, and \textbf{\textcolor{see}{SEE $\approx 4.54\times10^{11}$\,h/cm$^2$}}.
    \end{itemize}
    \vspace{0.12cm}
    \begin{center}
        \includegraphics[width=0.95\linewidth]{assets/fig_radiation_criteria.pdf}
    \end{center}
    \vspace{-0.2cm}
    \captiontext{Worked example for the LVPS bricks: simulated dose and fluence are scaled by safety factors to establish the final qualification target.}
}

\block{Irradiation Test Facilities}{
    \looseitems
        \item The testing campaigns are distributed across specialised European facilities to accurately reproduce the diverse HL-LHC radiation environment.
        \item Each facility targets a specific damage mechanism corresponding to \textcolor{tid}{TID}, \textcolor{niel}{NIEL}, or \textcolor{see}{SEE}, allowing for independent characterization of failure modes.
    \end{itemize}
    \vspace{0.5cm}
    \begin{center}
    \footnotesize
    \renewcommand{\arraystretch}{1.5}
    \begin{tabular}{p{0.18\linewidth} p{0.25\linewidth} p{0.50\linewidth}}
    \hline
    \textbf{Facility} & \textbf{Institution} & \textbf{Target Effect} \\
    \hline
    \textbf{CC60} & CERN & \textcolor{tid}{Total Ionizing Dose (TID)} via Gamma ($^{60}$Co) \\
    \textbf{TRIGA} & JSI (Ljubljana) & \textcolor{niel}{Non-Ionizing Energy Loss (NIEL)} via Reactor Neutrons \\
    \textbf{PIF} & PSI (Villigen) & \textcolor{see}{Single Event Effects (SEE/SEL)} via Protons \\
    \textbf{CHARM} & CERN & Mixed Field for System-Level Tests \\
    \hline
    \end{tabular}
    \end{center}
    \vspace{-0.2cm}
    \captiontext{Summary of primary irradiation facilities utilized for the TileCal qualification.}
}

\block{Optical Link (SFP+) Evaluation}{
    \looseitems
        \item Four different SFP+ variants were tested for the multi-Gbps optical links: AVAGO, CORETEK, FS-Europe, and Direktronik.
        \item The AVAGO (AFBR-709SMZ) and CORETEK devices failed the stringent irradiation tests early in the campaign.
        \item Conversely, both the Direktronik and FS-Europe transceivers demonstrated excellent radiation hardness, surviving fluences of $21\times 10^{12}$\,n/cm$^2$ and over 200\,Gy with no functional failures.
        \item The qualified transceivers ensure reliable data offloading at speeds of 9.6\,Gbps per link, providing substantial margin for the HL-LHC data rates.
    \end{itemize}
}

% ===========================================================================
%  COLUMN 2 
% ===========================================================================
\column{0.333}

\block{Daughterboard Test Campaigns}{
    \looseitems
        \item The Link and Control Daughterboard (DB) distributes clock and control signals while aggregating digitized data.
        \item Radiation campaigns involved irradiating fully assembled DB prototypes up to 108\,Gy and $14\times 10^{12}$\,n/cm$^2$.
        \item Key commercial components were qualified, including ProASIC FPGAs, IS25LP series flash memories, and specific oscillators.
    \end{itemize}
    \vspace{0.12cm}
    \begin{center}
        \includegraphics[width=0.90\linewidth]{assets/fig_db6_tid_test.png}
    \end{center}
    \vspace{-0.2cm}
    \captiontext{Power consumption of the DB6-2 board under continuous $^{60}$Co gamma irradiation, demonstrating stability up to 220\,Gy.}
}

\block{FPGA \& Memory Sensitivities}{
    \looseitems
        \item \textbf{Kintex Ultrascale (KU) FPGAs:} Initial tests on KU+ variants showed unacceptable \textbf{\textcolor{sel}{SEL}} rates (2\,SEL/$10^9$\,p/cm$^2$). Switching to the standard KU FPGA eliminated SEL up to $3.2\times 10^{11}$\,p/cm$^2$.
        \item The measured \textbf{\textcolor{see}{SEU}} rate in KU FPGAs translates to roughly 5498 upsets per year. Thanks to Xilinx Soft Error Management (SEM), only $\sim$11 uncorrectable errors are expected annually.
        \item \textbf{Flash Memories:} Devices like IS25LP128F exhibited temporary write-failures at high \textbf{\textcolor{tid}{TID}} but fully recovered re-programmability after room-temperature annealing.
    \end{itemize}
}

\block{Mitigation of Single Points of Failure}{
    \looseitems
        \item To ensure robust 10-year operation, the electronics architecture embeds multiple layers of active and passive redundancy.
        \item Data transmission utilizes rad-hard GBTx ASICs alongside the GBT-FEC (Forward Error Correction) protocol to safeguard against bit-flips in the optical links.
        \item Inside the FPGAs, Triple Mode Redundancy (TMR) protects mission-critical control logic from transient upsets, virtually eliminating the impact of radiation on nominal data taking.
        \item The Xilinx Soft Error Management (SEM) IP core continuously scrubs the FPGA configuration memory to detect and correct single-bit errors in the background.
        \item This multi-layered fault-tolerant approach ensures that temporary radiation effects do not permanently degrade detector performance.
    \end{itemize}
}

\block{Pre-Production \& Future Testing}{
    \looseitems
        \item Most components have been acquired and qualified on a lot-by-lot basis for both pre-production and full production phases.
        \item Final radiation campaigns to test the newly acquired lots of DC-DC converters (LTM4691), ProASIC FPGAs, and KU FPGAs are scheduled for 2025 and 2026.
        \item A complete system-level test of the final production DB and LVPS bricks in a mixed-field environment will ensure no latent weaknesses remain.
        \item These forthcoming campaigns will validate the minor layout modifications introduced to accommodate the final component footprints.
        \item A rigorous quality assurance pipeline will be maintained throughout the entire 2026--2027 manufacturing period.
    \end{itemize}
}

\block{LVPS System-Level Qualification}{
    \looseitems
        \item The Low-Voltage Power Supply (LVPS) employs rad-tolerant bricks to step down 200\,VDC to 10\,VDC.
        \item System-level testing at the CHARM mixed-field facility exposed unexpected failures not seen in individual component tests.
        \item Power MOSFETs experienced a gate-source threshold shift to zero at 323\,Gy, causing a short-circuit, though this heavily exceeds the 80\,Gy requirement.
    \end{itemize}
    \vspace{0.12cm}
    \begin{center}
        \includegraphics[width=0.90\linewidth]{assets/fig_lvps_tests_overview.pdf}
    \end{center}
    \vspace{-0.2cm}
    \captiontext{Overview of the LVPS test setup inside the CHARM mixed-field facility, used for evaluating combined radiation effects.}
}

% ===========================================================================
%  COLUMN 3 
% ===========================================================================
\column{0.333}

\block{LVPS Brick Architecture}{
    \looseitems
        \item The LVPS relies on radiation-tolerant "bricks" designed to step down the external 200\,VDC bulk power to 10\,VDC for the on-detector Front-End electronics.
        \item These modular power supplies must operate within the restricted and highly irradiated volume of the TileCal structural fingers.
    \end{itemize}
    \vspace{0.12cm}
    \begin{center}
        \includegraphics[width=0.92\linewidth]{assets/fig_lvps_brick.pdf}
    \end{center}
    \vspace{-0.2cm}
    \captiontext{Physical architecture of the upgraded LVPS brick highlighting the dense integration of rad-tolerant components.}
}

\block{Cable \& Routing Optimization}{
    \looseitems
        \item High radiation zones exacerbate cable degradation. To address this, the 200\,VDC power routing was completely overhauled for the HL-LHC era.
        \item The number of main cables per module was reduced from four to two, significantly shortening the total exposed length in the highest radiation regions.
        \item This routing optimization resulted in a \textbf{94\% reduction} in the accumulated dose on the vulnerable segments of the cables entering the LVPS bricks.
    \end{itemize}
    \vspace{0.12cm}
    \begin{center}
        \includegraphics[width=0.92\linewidth]{assets/fig_wire_reduction.pdf}
    \end{center}
    \vspace{-0.2cm}
    \captiontext{Comparison of legacy versus optimized 200\,V routing schemes, minimizing radiation exposure.}
}

\block{LVPS Component Selection \& Evaluation}{
    \looseitems
        \item A rigorous campaign tested commercial-off-the-shelf (COTS) components for the Low-Voltage Power Supply (LVPS).
        \item Key components were irradiated individually prior to full system tests.
    \end{itemize}
    \vspace{0.5cm}
    \begin{center}
    \footnotesize
    \renewcommand{\arraystretch}{1.5}
    \begin{tabular}{p{0.22\linewidth} p{0.18\linewidth} p{0.52\linewidth}}
    \hline
    \textbf{Component} & \textbf{Function} & \textbf{Radiation Test Results} \\
    \hline
    \textbf{LT1681} & Controller & Passed \textcolor{tid}{TID}, \textcolor{niel}{NIEL}, \textcolor{see}{SEE} up to target. \\
    \textbf{LTC6241} & Op-Amp & Passed. Minor SETs ($\sim$50\,mV) observed. \\
    \textbf{LT3080} & Voltage Reg. & Passed CHARM (835\,Gy, $6.57\times 10^{12}$\,n/cm$^2$). \\
    \textbf{IR2110} & MOSFET Drv. & Passed. No SEU up to $3.4\times 10^{11}$\,p/cm$^2$. \\
    \textbf{SIHFS9N60A} & Power MOSFET & Passed. Threshold shift at 200\,Gy. \\
    \textbf{SI8920} & Isolation Amp. & Passed \textcolor{tid}{TID}/\textcolor{see}{SEE} for input $<100$\,mV. \\
    \textbf{HCPL7800} & Isolation Amp. & Failed \textcolor{niel}{NIEL} at $2.75\times 10^{12}$\,n/cm$^2$. \\
    \hline
    \end{tabular}
    \end{center}
    \vspace{-0.2cm}
    \captiontext{Overview of critical LVPS component qualification outcomes.}
}

\block{MOSFET Threshold Shifts}{
    \looseitems
        \item Ionizing radiation alters trapped charge in MOSFET oxides, systematically shifting the gate threshold voltage.
        \item FDFMA3N109 MOSFETs used for digital switching were monitored during \textbf{\textcolor{tid}{TID}} testing.
        \item Despite a progressive threshold decrease, post-irradiation values remained well within digital logic margins.
    \end{itemize}
    \vspace{0.12cm}
    \begin{center}
        \includegraphics[width=0.90\linewidth]{assets/fig_mosfet_threshold.png}
    \end{center}
    \vspace{-0.2cm}
    \captiontext{Shift in the gate threshold voltage of the FDFMA3N109 MOSFET.}
}

\block{Conclusions \& Outlook}{
    \looseitems
        \item A comprehensive radiation methodology using independent \textcolor{tid}{TID}, \textcolor{niel}{NIEL}, and \textcolor{see}{SEE} facilities, augmented by mixed-field testing, is essential for HL-LHC electronics qualification.
        \item The DB design is fully qualified. Over 930 boards will be produced for the upgrade.
        \item The LVPS brick design has been optimized based on CHARM results, substituting the HCPL7800 with the SI8920 amplifier.
        \item Pre-production batches will continue to undergo lot-acceptance testing following ATLAS guidelines before final deployment.
    \end{itemize}
}

\end{columns}
\end{document}
